\documentclass[10pt,a4paper]{amsart}
\usepackage[margin=19mm]{geometry}
\usepackage{amsmath,amssymb,mathtools}
\usepackage[scr=rsfs,cal=euler]{mathalfa}
\usepackage{xfrac}
\usepackage[unicode,hidelinks,bookmarks=false]{hyperref}
\newcommand{\ie}{i.e.\ }
\newcommand{\cf}{cf.\ }
\newcommand{\resp}{respectively\ }
\newcommand{\Subsection}{\subsection}
\providecommand{\nobreakdash}{\nobreak\hbox{-}}
\setlength{\emergencystretch}{2em}
\multlinegap0pt
\usepackage{tikz-cd}
\usepackage[all]{xy}
\let\scr\mathscr % source script-letter compatibility via installed mathalfa
\newtheorem{theorem}{Theorem}
\newtheorem{lemma}{Lemma}
\theoremstyle{definition}
\newtheorem{definition}{Definition}
\newtheorem{remark}{Remark}
 \newcommand{\Wi}{\widetilde}
\DeclareMathOperator{\cat}{{\mathsf{cat}}}
\DeclareMathOperator{\secat}{{\mathsf{secat}}}
\DeclareMathOperator{\dsecat}{{\mathsf{dsecat}}}
\DeclareMathOperator{\dcat}{{\mathsf{dcat}}}
\def\B{{\scr B}}
\title{Multiplicative bound for distributional category along a covering}
\author{Ekansh Jauhari, John Oprea}
\date{Source: arXiv:2508.06973v2 (CC BY 4.0)}
\begin{document}
\maketitle
\noindent\emph{Source and licence.} On distributional one-category, \\ diagonal distributional complexity, \\ and related invariants. arXiv:2508.06973v2 (CC BY 4.0). The official fixed-version abstract page explicitly links to Creative Commons Attribution 4.0 International. The saved HTML and response record were checked on 2026-10-10. See \url{https://arxiv.org/abs/2508.06973v2} and \url{https://creativecommons.org/licenses/by/4.0/}. The licence witness is bundled in \texttt{licenses/license-evidence-182976.json}; no original publisher-PDF inspection is claimed.\par\smallskip
\noindent\textit{Editorial scope.} The single counted unit is the theorem labelled \texttt{thm:general ineq} (source lines 1827--1832), with its entire original proof (1833--1884). The uncounted preliminary passages supply the connected-CW convention, based-path evaluation, finite-support probability spaces and their fiberwise naturality, the definitions of distributional category and one-category, the pullback lemma, and the explicit universal-cover lifting construction. All retained mathematical passages are copied verbatim. The sole reference adaptation changes the motivational pointer \texttt{cor:isoineq} at source line 799 to a link identifying that label in the fixed-version full article; that result is not used as a premise of the counted proof. The continuation through source line 861 is retained to close the universal-cover paragraph and its displayed definitions. Original editable commutative diagrams and the xymatrix are retained; no external image file is added. The cited bibliography entries are transcribed from the same source, with their source labels. Wrapper section and theorem numbering is editorial; the source script-letter command uses the installed RSFS font through mathalfa in place of the unavailable rsfso package. The full native \texttt{sources/182976/source0.tex} is unchanged. No new mathematical proof or correction is supplied.
\section*{Uncounted source-local prerequisites}
\subsection*{Notations and conventions} All topological spaces considered in this paper are connected CW complexes. We use the term \emph{maps} for continuous functions. The symbol ``$\cong$" is used to denote homeomorphisms, isomorphisms, and bijections, and the symbol ``$\simeq$" is used to denote homotopy equivalences of maps and spaces. 

\subsection{Sectional category}\label{subsec:secat section}
Suppose $f:E\to B$ is a (Hurewicz) fibration. For each $n \ge 1$, the iterated join of $n$ copies of $E$ along $f$ is
\[
\ast^n E = \left\{\sum_{i=1}^{n} \lambda_{i} e_{i} \hspace{1mm} \middle| \hspace{1mm} e_{i} \in E, \hspace{0.5mm} \sum_{i=1}^{n}\lambda_{i} = 1, \hspace{0.5mm} \lambda_{i} \geq 0, \hspace{0.5mm} f(e_{i})= f(e_{j})\right\},
\]
and the iterated fibration $\ast^nf : \ast^nE \to B$ is defined as
\[
\ast^n f \left(\sum_{i=1}^{n} \lambda_{i} e_{i}\right) = f(e_i)
\]
for any $i$ with $\lambda_{i} > 0$. The smallest integer $n$ such that $\ast^{n+1}f$ has a section is called the \emph{sectional category} of $f$, denoted $\secat(f)$, see~\cite{Sch}. It is well-known that $\secat$ is a homotopy invariant of fibrations.

\begin{remark}\label{rmk:homotopy section}
    If $f:E\to B$ is a fibration that admits a homotopy section, then $f$ admits a true section. This follows from the homotopy lifting property of $f$, see, for example,~\cite[Exercise 1.22]{CLOT}. So, in this paper, we will use the term ``section" for both a homotopy section and a true section of a given fibration since one can be obtained from the other. Similarly, since a lift of a map $g:A\to B$ along a fibration $f$ is essentially a section of the pullback fibration $g^*f$, we will use the term ``lift" for both a homotopy lift as well as a true lift of a given fibration along a given map.
\end{remark}

If $E=P_0(X)$ is the based path space (containing paths \emph{terminating} at a fixed basepoint $x_0 \in X$) and $f$ is the evaluation-at-zero map, defined $e^X(\phi)=\phi(0)$, then the least $n$ such that $\ast^{n+1}e^X$ admits a section is called the \emph{Lusternik--Schnirelmann category} of $X$, denoted $\cat(X)$. More generally, for a map $f:Y\to X$, the least $n$ such that $f$ has a lift along $\ast^{n+1}e^X$ is called the \emph{Lusternik--Schnirelmann category} of $f$, denoted $\cat(f)$. Clearly, $\cat(X)=\cat(1_X)$. Since $\cat$ is a homotopy invariant and the
homotopy type of an Eilenberg--Mac~Lane space $K(\pi,1)$ is determined by $\pi$, we 
may use the notation $\cat(K(\pi,1)) = \cat(\pi)$.

\subsection{Distributional sectional category}\label{sec:basic2}

Let $\B(Z)$ denote the set of probability measures on a topological space $Z$ and 
$$\B_n(Z) = \{\mu \in B(Z) \mid\, |\text{supp}(\mu)| \leq n\}$$ 
be the subspace of measures supported by at most $n$ points. So, each such measure $\mu\in\B_n(Z)$ 
can be written as a formal sum 
$$\mu = \sum_{i=1}^n a_i \delta_{z_i} = \sum_{i=1}^{n} a_i z_i,$$
where $z_i \in Z$, $\delta_{z_i}$ is the point (Dirac) measure at $z_i$, and 
$a_i \geq 0$ with $\sum_{i=1}^n a_i = 1$. We identify the space $Z$ with a 
subspace of $\B_n(Z)$ by means of Dirac measures $\delta_z$ for every $z \in Z$.
The support of $\mu$ is the subset $\text{supp}(\mu) = \{z_i\in Z\mid\, a_i > 0\}$. In the case when $Z$ is a metric space, we consider 
the L\'evy-Prokhorov metric on $\B_n(Z)$, see~\cite{Pr,DJ}. A different topology is defined 
in~\cite{KK,KW}, but for the spaces we consider, the results of~\cite{KW} also hold, since~\cite{KW2} shows that the distributional invariants agree with their analog counterparts on compact metric spaces and finite groups, among other spaces.

If $f\colon E \to B$ is a map, then for each integer $n\ge 1$, there is a space
\[
E_n = \bigcup_{b\in B}\B_n\left(f^{-1}(b)\right)=\left\{\mu \in \B_n(E)\mid\, \text{supp}(\mu) \subset f^{-1}(b)\ \text{for\ some\ } b \in B\right\},
\]
which is topologized as a subspace of $\B_n(E)$, and a map $f_n\colon E_n \to B$ defined as
\[
f_n(\mu) = b \ \text{ if } \mu\in\B_n\left(f^{-1}(b)\right).
\]


\begin{remark}[\protect{Naturality}]\label{rem:naturality}
    Suppose we have the following commutative diagram:
\[
\begin{tikzcd}
    X \arrow{r}{f}  \arrow{d}{p}
    &
    Y \arrow{d}{q}
    \\
    B \arrow{r}{g}
    &
    C.
\end{tikzcd}
\]
Of course, for any integer $n\ge 1$, we have maps $p_n\colon X_n\to B$ and $q_n\colon Y_n\to C$ as defined above. Applying the functor $\B_n$ on $f$ yields a map $\B_n(f)\colon \B_n(X)\to \B_n(Y)$. Let $f_n$ be the restriction of $\B_n(f)$ to the space $X_n\subset \B_n(X)$. It is easy to check that the image of $f_n$ lies in $Y_n\subset \B_n(Y)$, and $q_n\circ f_n=g\circ p_n$ --- see, for example, the proof of~\cite[Proposition~5.3]{Ja1}. Hence, given a map $f\colon X\to Y$ over $g\colon B\to C$, we get a natural map $f_n\colon X_n\to Y_n$ also over $g\colon B\to C$. This fact will be used frequently in our proofs.
\end{remark}

If $f\colon E\to B$ is a fibration with fiber $F$, then $f_n\colon E_n \to B$ is a fibration with fiber $\B_n(F)$, see~\cite{DJ,KW}. If $F$ is $r$-connected, then $\B_n(F)$ is $(2n+r-2)$-connected for $r\ge 1$ (see~\cite{KK}), $(n-1)$-connected for $r=0$, and $(n-2)$-connected otherwise, see~\cite[Section 3]{Dr}. The smallest integer $n$ such that $f_{n+1}\colon E_{n+1}\to B$ has a section is called the \emph{distributional sectional category} of $f$, denoted $\dsecat(f)$. If two fibrations are fiberwise homotopic, then their $\dsecat$ values coincide,~\cite{Ja1}. We refer to~\cite[Section 2.1]{Ja2} for a motivation behind this invariant. The similar invariant defined in~\cite{KW} is called 
``analog sectional category" and denoted by $\sf{{asecat}}$$(f)$. 

If $E = P_0(X)$ and $f=e^X$, then we write
$e_n^X \colon P_0(X)_n \to X$, and the least $n$ such that $e_{n+1}^X$ has a section is called the \emph{distributional category} of $X$, denoted $\dcat(X)$, see~\cite{DJ} (a similar invariant was denoted $\sf{{acat}}$$(X)$ in~\cite{KW}). If $f\colon Y \to X$
is a map, then the least $n$ such that $f$ has a lift $\widetilde f$ in
\[
\begin{tikzcd}[contains/.style = {draw=none,"\in" description,sloped}]
&
P_0(X)_{n+1} \arrow{d}{e_{n+1}^X}
\\
Y  \arrow{r}{f} \arrow{ur}{\widetilde{f}}
&
X
\end{tikzcd}
\]
is called the \emph{distributional category} of $f$, denoted $\dcat(f)$, see~\cite{Dr,Ja2}. From the proof of Lemma~\ref{lem:dsecat1} below, one can see that $\dcat(f)=\dsecat(f^*e^X)$, where $f^*e^X$ is the pullback of the evaluation fibration $e^X$ along $f$. In particular, we have that $\dcat(X)=\dcat(1_X)=\dsecat(e^X)$.
We note that this definition is equivalent to the original one in~\cite{DJ}, which fits better with the viewpoint of the Introduction.

\begin{lemma}\label{lem:dsecat1}
    If $p:E\to B$ is a fibration and $f:A\to B$ is a map, then $\dsecat(f^*p)\le\dsecat(p)$, where $f^*p$ is the pullback of $p$ along $f$.
\end{lemma}

\begin{proof}
    Let $\dsecat(p)=n$ and $s:B\to E_{n+1}$ be a section of $p_{n+1}:E_{n+1}\to B$. For each $x\in A$, we have a homeomorphism between $(f^*p)^{-1}(x)$ and $p^{-1}(f(x))$.
    %The fiber of $x\in A$ under $f^*p$ is homeomorphic to the fiber of $f(x)\in B$ under $p$.
    Since the functor $\B_{n+1}$ preserves homeomorphisms, the fibers of the fibrations $p_{n+1}$ and $(f^*p)_{n+1}:(f^*E)_{n+1}\to A$ are homeomorphic and we have the homeomorphism
    \[
    f^*E_{n+1}=\bigcup_{x\in A} \B_{n+1}\left(p^{-1}(f(x))\right)=\bigcup_{x\in A} \B_{n+1}\left((f^*p)^{-1}(x)\right)=(f^*E)_{n+1}.
    \]
    Hence, we have the following (homotopy) pullback diagram:
    \[\begin{tikzcd}
     A\arrow[ddr, bend right, "1_B"] \arrow[drr, bend left, "s\circ f"] \arrow[dr, dashed, "\tau"] \\
 &
    (f^*E)_{n+1} \arrow{r} \arrow{d}{(f^*p)_{n+1}} & E_{n+1} \arrow[d, "p_{n+1}"] \\
    & A \arrow[swap]{r}{f}                  & B,
    \end{tikzcd}\]
    where $\tau:A\to (f^*E)_{n+1}$ exists by the universal property of the pullback. Thus, we get a section of $(f^*p)_{n+1}$ and hence the inequality $\dsecat(f^*p)\le n$.
\end{proof}

\begin{definition}
    If $q\colon \Wi{X}\to X$ is the universal covering, then the least integer $n$ such that $q_{n+1}^X$ has a section is called the \emph{distributional one-category} of $X$, denoted $\dcat_1(X)$. More generally, if
$f\colon Y \to X$ is a map, then the least $n$ such that $f$ has a lift $\widetilde f$ in
\[
\begin{tikzcd}[contains/.style = {draw=none,"\in" description,sloped}]
&
\widetilde X_{n+1} \arrow{d}{q_{n+1}^X}
\\
Y  \arrow{r}{f} \arrow{ur}{\widetilde{f}}
&
X
\end{tikzcd}
\]
is called the \emph{distributional one-category} of $f$, denoted $\dcat_1(f)$.
\end{definition}

\subsection{Behavior on covering maps}\label{subsec: behavior covering}
We now aim to generalize Corollary~\href{https://arxiv.org/abs/2508.06973v2}{cor:isoineq} from isomorphisms to monomorphisms. To do that, we first look at covering spaces. In this paper, the symbol $\simeq_*$ means homotopy rel endpoints for paths.

\begin{lemma}\label{lem:lifthomo}
Suppose $p \colon Y \to X$ is a covering and $\omega$, $\sigma$ are two paths in $Y$ with $\omega(0)
= \sigma(0)$ and $p\circ \omega \simeq_* p \circ \sigma$. Then $\omega \simeq_* \sigma$.
\end{lemma}
\begin{proof}
Let $H\colon I \times I\to X$ be a homotopy with $H_0=p\circ \omega$ and $H_1 = p\circ \sigma$ such that
$H(0,s)= p(\omega(0)) = p(\sigma(0))$ and $H(1,s) = p(\omega(1)) = p(\sigma(1))$ for each $s\in I$. Here, we use the convention $H_s(t)=H(t,s)$. We have a lift
$G$ in the commutative diagram
\[
\begin{tikzcd}[contains/.style = {draw=none,"\in" description,sloped}]
I \times \{0\} \arrow{r}{\omega} \arrow[swap]{d}
&
Y \arrow{d}{p}
\\
I\times I  \arrow{r}{H} \arrow[ur, dashed, "G"]
&
X.
\end{tikzcd}
\]
Therefore, $p\circ G=H$ with $G_0=\omega$ and $p\circ G_1=H_1=p\circ \sigma$. We thus have $p(G(0,s))=H(0,s)=p(\omega(0))
=p(\sigma(0))$ and $p(G(1,s))=H(1,s)=p(\omega(1))=p(\sigma(1))$. Hence, $G(0,s) \in p^{-1}(p(\omega(0)))$
and $G(1,s) \in p^{-1}(p(\omega(1)))$. These fibers are discrete and $G$ is continuous, so for all $s\in I$, we must have
$G(0,s)=\bar y$ and $G(1,s)=\widehat y$ for some points $\bar y,\widehat y\in Y$. Now, $G(0,0)=G_0(0)=\omega(0)$, so $\bar y=\omega(0)$. Also, $G(1,0)=G_0(1)=\omega(1)$, so $\widehat y=\omega (1)$. Note that $G_1(0)=G(0,1)=\bar y=\omega(0)=\sigma(0)$ and $p\circ G_1=p\circ \sigma$. By the unique path lifting property of $p$, we must have $G_1 = \sigma$. Therefore, $G$ gives
a homotopy $\omega =G_0\simeq_* G_1 = \sigma$. Observe that $\sigma(1)=G_1(1)=\widehat y=\omega(1)$, so $\omega$ and $\sigma$ have the same terminal points.
\end{proof}

\begin{remark}\label{rem:lifthomo2}
    It follows from the technique of the proof of Lemma~\ref{lem:lifthomo} that $\omega\simeq_{*}\sigma$ is true again if $\omega(1)=\sigma(1)$ and $p\circ \omega\simeq_{*}p\circ\sigma$.
\end{remark}

Given a covering map $p:Y\to X$, we use the definition of universal coverings to get
\[
\widetilde X = P_0(X)/[-]\ \text{ and }\ \widetilde Y = P_0(Y)/\{-\},
\]
where $[\alpha]=[\beta]$ for $\alpha,\beta\in P_0(X)$ if $\alpha(0)=\beta(0)$ and $\alpha\simeq_* \beta$, and $\{\sigma\}
=\{\tau\}$ for $\sigma,\tau \in P_0(Y)$ if $\sigma(0)=\tau(0)$ and $\sigma\simeq_* \tau$. We recall that $P_0(X)$ consists of paths in $X$ having terminal point $x_0$, and $P_0(Y)$ consists of paths in $Y$ having terminal point $y_0\in p^{-1}(x_0)$. The universal covering maps are given by $q^X([\gamma])=\gamma(0)$ and $q^Y(\{\alpha\})=\alpha(0)$. Of course, we know that $\widetilde X \cong \widetilde Y$, but we need the exact way this happens.

Define $P\colon \widetilde Y \to \widetilde X$ by $P(\{\sigma\}) = [p\circ\sigma]$. Clearly if
$\sigma\simeq_*\tau$, then $p\circ\sigma\simeq_* p\circ\tau$, so $P$ is well-defined.
Define $\theta\colon \widetilde X \to \widetilde Y$ as follows. Let $[\gamma]\in \widetilde X$. Because $\gamma(1)=x_0$, there is a unique lift $\widetilde\gamma\colon I \to Y$ to $\gamma$ ending at
$\widetilde\gamma(1)=y_0$. Define $\theta([\gamma])=\{\widetilde\gamma\}$. This map is well-defined
by Remark~\ref{rem:lifthomo2} since $\gamma \simeq_* \delta$ says that the unique 
lifts $\widetilde\gamma$
and $\widetilde \delta$ also have $\widetilde\gamma \simeq_* \widetilde\delta$ because $\widetilde\gamma(1) = y_0 =
\widetilde\delta(1)$. It is easy to check that
\[
P\circ\theta([\gamma])  = P(\{\widetilde\gamma\}) = [p\circ\widetilde\gamma]  = [\gamma] \ \text{ and } \
\theta\circ P(\{\sigma\}) = \theta([p\circ\sigma]) = \{\sigma\},
\]
where the last equality comes from the unique path lifting property of $p$. Now, recall that 
\[
\widetilde X_{n+1} = \{\mu \in \B_{n+1}(\widetilde X) \mid\, \text{supp}(\mu) \subset (q^X)^{-1}(x)\
\text{for\ some\ } x\in X\}, \ \text{ and}
\]
\[
\widetilde Y_{n+1} = \{\nu \in \B_{n+1}(\widetilde Y) \mid\, \text{supp}(\nu) \subset (q^Y)^{-1}(y)\
\text{for\ some\ } y\in Y\}.
\]
So, $\mu = \sum_i a_i [\gamma_i] \in \widetilde{X}_{n+1}$ means there exists $x\in X$ such that $q^X(\gamma_i)=\gamma_i(0)=x$ is independent of $i$, and $\nu=\sum_i b_i \{\alpha_i\} \in \widetilde{Y}_{n+1}$ means there exists $y\in Y$ such that $q^Y(\alpha_i)=\alpha_i(0)=y$ is independent of $i$.

\section*{Counted theorem and complete author proof}
\begin{theorem}\label{thm:general ineq}
    If $p\colon Y\to X$ is a covering map, then
    \[
    \dcat(X)\le (\dcat_1(X)+1)(\dcat(Y)+1)-1.
    \]
\end{theorem}

\begin{proof}
    Let $x_0\in X$ be the basepoint of $X$ and $y_0\in p^{-1}(x_0)$ be the basepoint of $Y$. Recall from Section~\ref{subsec: behavior covering} that $\Wi{X}=P_0(X)/[-]$ and $\Wi{Y}=P_0(Y)/\{-\}$, and $\theta\colon \Wi X\to\Wi Y$ is a homeomorphism which maps any $[\gamma]\in \Wi X$ to $\{\Wi{\gamma}\}\in \Wi Y$, where $\Wi{\gamma}\in P_0(Y)$ is the lift of $\gamma\in P_0(X)$ with respect to the point $\gamma(1)=x_0$. Recall also that the universal coverings $q^X\colon \Wi X\to X$ and $q^Y\colon \Wi Y\to Y$ satisfy $q^X[\alpha]=\alpha(0)$ and $q^Y\{\beta\}=\beta(0)$. Now, the homeomorphism $\theta$ 
satisfies $p\circ q^Y \circ \theta = q^X$, so we have a diagram
$$\xymatrix{
\Wi X \ar[rr]^-\theta \ar[dr]_-{q^X} & & \Wi Y \ar[dl]^-{p\circ q^Y} \\
& X &
} 
$$
that induces a map $\theta_{n+1}\colon \Wi X_{n+1} \to \Wi Y^p_{n+1}$, where $\Wi Y^p_{n+1}$
denotes the distributional construction on the covering $p \circ q^Y$. This means that an element
\[
\sum_{i=1}^{n+1} b_i \Wi{y}_i  \in \Wi Y^p_{n+1}
\]
has all $\Wi{y_i}$ in the same fiber of $p \circ q^Y$, but not necessarily in the same fiber
of $q^Y$. Essentially, we are taking points in the fiber of $p$ and then taking points in 
the fibers over these points. 
This explains the notation $\Wi Y^p_{n+1}$. Also, the map $q^Y$ induces a map $(q^Y_p)_{n+1} 
\colon \Wi Y^p_{n+1} \to Y_{n+1}$, where $Y_{n+1}$ is, as usual, the distributional construction applied to the covering $p\colon Y \to X$. 

Now, let $\dcat_1(X)=n$ and $\dcat(Y)=r$. We then have a section $\sigma\colon X\to \Wi X_{n+1}$ of $q^X_{n+1}$. For $x\in X$, suppose
    \[
    \sigma(x)=\sum_{i=1}^{n+1}a_i\Wi{x}_i \ \text{ with } \ q^X(\Wi{x}_i)=x.
    \]
    For each $1\le i \le n+1$, we can write $\Wi{x}_i=[\gamma_i]$ for some $\gamma_i\in P_0(X)$ 
such that $\gamma_i(0)=x$. Then, for $\theta_{n+1}\circ \sigma\colon X\to \Wi Y^p_{n+1}$, we have that
    \[
    \theta_{n+1}\circ \sigma(x) =\sum_{i=1}^{n+1}a_i\{\Wi{\gamma}_i\} \ \text{ with } \ p(\Wi{\gamma_i}(0))=x.
    \]
(Note that this element is not in $\Wi Y_{n+1}$ since the $\Wi{\gamma_i}$ are not
necessarily in the same fiber of $q^Y$.) 
For simplicity, let us write $\Wi{\gamma_i}(0)=y_i\in Y$. Then,
    \[
    (q^Y_p)_{n+1}\circ \theta_{n+1}\circ \sigma(x) =\sum_{i=1}^{n+1}a_iy_i \ \text{ with } \ p(y_i)=x.
    \]
    Since $\dcat(Y)=r$, the map $e^Y_{r+1}$ has a section $s\colon Y\to P_0(Y)_{r+1}$. 
For each $y_i$ as above, suppose
\[     
s(y_i)=\sum_{j=1}^{r+1}b_j\alpha_{ij} \ \text{ with } \ \alpha_{ij}(0)=y_i.       
\]
    Let $k=(n+1)(r+1)$. Suppose $\tau\colon P_0(Y)\to P_0(X)$ is the natural map defined 
as $\tau(\phi)=p\circ\phi$. This extends to a map $\tau_k\colon P_0(Y)_k \to P_0(X)_k$
since $p \circ e^Y = e^X \circ \tau$. Note that $\tau_k$ restricted to $P_0(Y)_{r+1}$
is $\tau_{r+1}$. Define a map $t\colon Y_{n+1}\to P_0(X)_k$ by
\[
t \left( \sum_{i=1}^{n+1}\lambda_iz_i \right) = \sum_{i=1}^{n+1}\lambda_i \tau_k\left(s(z_i)\right)
\]
for $z_i\in Y$. Then, feeding the measure $(q^Y_p)_{n+1}\circ \theta_{n+1}\circ \sigma(x) \in Y_{n+1}$ into $t$ yields
\[
t \left( \sum_{i=1}^{n+1}a_iy_i \right) = \sum_{i=1}^{n+1}a_i \tau_k\left(s(y_i)\right) = \sum_{i=1}^{n+1}\sum_{j=1}^{r+1}a_ib_j\left(p\circ\alpha_{ij}\right).
\]
Note that $e^Y(\alpha_{ij})=y_i$ for all $j$ and $e^X(p\circ\alpha_{ij})=x$ for all $i$ and $j$, So, the composition $t\circ (q^Y_p)_{n+1}\circ \theta_{n+1}\circ \sigma\colon X\to P_0(X)_k$ is a section of the fibration $e_k^X\colon P_0(X)_k\to X$. Therefore, $\dcat(X) \leq k-1 = (n+1)(r+1)-1$.
\end{proof}

\begin{thebibliography}{99}
\bibitem[CLOT03]{CLOT} O.~Cornea, G.~Lupton, J.~Oprea, D.~Tanré, \emph{Lusternik--Schnirelmann Category}. Math. Surveys Monogr., 103, AMS, Providence, 2003, xviii+330 pp.
\bibitem[Dr]{Dr} A.~Dranishnikov, Distributional topological complexity of groups. Preprint, arXiv:2404.03041 [math.GT] (2024), pp.~29.
\bibitem[DJ24]{DJ} A.~Dranishnikov, E.~Jauhari, Distributional topological complexity and LS-category. In \emph{Topology and AI---topological aspects of algorithms for autonomous motion}, ed. M.~Farber and J.~Gonz\'alez, EMS Ser. Ind. Appl. Math., 4, EMS Press, Berlin, 2024, pp. 363–385.
\bibitem[Ja25a]{Ja1} E.~Jauhari, On sequential versions of distributional topological complexity. \emph{Topology Appl.} \textbf{363} (2025), 109271, pp.~28.
\bibitem[Ja25b]{Ja2} E.~Jauhari, Distributional category of manifolds. \emph{Bol. Soc. Mat. Mex.} \textbf{31} (2025), no.~2, 63, pp.~34.
\bibitem[KK11]{KK} S.~Kallel, R.~Karoui, Symmetric joins and weighted barycenters. \emph{Adv. Nonlinear Stud.} \textbf{11} (2011), no.~1, 117--143.
\bibitem[KW24]{KW} B.~Knudsen, S.~Weinberger, Analog category and complexity. \emph{SIAM J. Appl. Algebra Geom.} \textbf{8} (2024), no.~3, 713--732.
\bibitem[KW25]{KW2} B.~Knudsen, S.~Weinberger, On the analog category of finite groups. To appear in \emph{Algebr. Geom. Topol.}, arXiv:2411.14958 [math.AT], pp. 11.
\bibitem[Pr56]{Pr} Y.~V.~Prokhorov, Convergence of random processes and limit theorems in probability theory. \emph{Theory Probab. Appl.} \textbf{1} (1956), 157--214. 
\bibitem[Sc66]{Sch} A.~S.~Schwarz, The genus of a fiber space. In \emph{Eleven papers on topology and algebra}, Amer. Math. Soc. Transl. Ser. 2, 55, AMS, Providence, 1966, pp.~49--140. 
\end{thebibliography}
\end{document}
